\begin{progprob}
    % Write the problem here.
    In a file named \python{swap_sum.py}, write a function named
    \python{swap_sum(A, B)} which, given two \textbf{sorted} integer arrays
    $\python{A}$ and $\python{B}$, returns a pair of indices \python{(A_i,
    B_i)} -- one from $A$ and one from $B$ -- such that after swapping these
    indices, \python{sum(B) == sum(A) + 10}. If more than one pair is found,
    return any one of them. If such a pair does not exist, return
    $\python{None}$. 

    For example, suppose \python{A = [1, 6, 50]} and \python{B = [4, 24, 35]}.
    Swapping $6$ and $4$ results in arrays $(1, 4, 50)$ and $(6, 24, 35)$; the
    elements of each list sum to 55 and 65. Thus, you must return $\python{(1,
    0)}$ as you are expected to return the indices. 

    Your algorithm should run in time $\Theta(n)$, where $n$ is the size of the
    larger of the two lists. Your code should not modify \python{A} and \python{B}.

    \emph{Hint}: This is similar to the movie problem from lecture, but also to
    a problem covered in discussion. In the discussion problem,
    we're given two lists, \python{A} and \python{B} and a target \python{t},
    and our goal is to find an element $a$ of \python{A} and an element $b$ of
    \python{B} such that $a + b = t$. This problem is similar, in that we want
    $(\text{sum of A after swap}) - (\text{sum of B after swap}) = -10$. Note
    that here we are subtracting, where in discussion we were adding! How does
    that change things?

    This is a coding problem, and you'll submit your \python{swap_sum.py}
    file to the Gradescope autograder assignment named ``Homework 03 -
    Programming Problem 02''. The public autograder will test to make sure your
    code runs without error on a simple test, so be sure to check its output!
    After the deadline a more thorough set of tests will be used to grade your
    submission.

    \begin{soln}

        \inputminted{python}{\thisdir/include-solution/swap_sum.py}

    \end{soln}
\end{progprob}
